\let\negmedspace\undefined
\let\negthickspace\undefined
\documentclass[journal]{IEEEtran}
\usepackage[a5paper, margin=10mm, onecolumn]{geometry}
\usepackage{lmodern} % Ensure lmodern is loaded for pdflatex
\usepackage{tfrupee} % Include tfrupee package

\setlength{\headheight}{1cm} % Set the height of the header box
\setlength{\headsep}{0mm}     % Set the distance between the header box and the top of the text

\usepackage{gvv-book}
\usepackage{gvv}
\usepackage{cite}
\usepackage{amsmath,amssymb,amsfonts,amsthm}
\usepackage{algorithmic}
\usepackage{graphicx}
\graphicspath{{./figs/}}
\usepackage{textcomp}
\usepackage{xcolor}
\usepackage{txfonts}
\usepackage{listings}
\usepackage{enumitem}
\usepackage{mathtools}
\usepackage{gensymb}
\usepackage{comment}
\usepackage[breaklinks=true]{hyperref}
\usepackage{tkz-euclide} 
\usepackage{listings}
\usepackage{gvv}                                        
\def\inputGnumericTable{}                    

\usepackage[latin1]{inputenc}                                
\usepackage{color}                                            
\usepackage{array}                        

\usepackage{longtable}                                       
\usepackage{calc}                            
\usepackage{multirow}             

\usepackage{hhline}                                          
\usepackage{ifthen}                              

\usepackage{lscape}
\usepackage{circuitikz}

\begin{document}
	
	\bibliographystyle{IEEEtran}
	\vspace{3cm}
	
	\title{12.263}
	\author{EE25BTECH11042 - Nipun Dasari}
	\maketitle
	
	\renewcommand{\thefigure}{\theenumi}
	\renewcommand{\thetable}{\theenumi}
	\setlength{\intextsep}{10pt}
	
	\numberwithin{equation}{enumi}
	\numberwithin{figure}{enumi}
	\renewcommand{\thetable}{\theenumi}
	
	\textbf{Question}:\\
	The matrix
	\begin{align}
		\vec{A} = \begin{myvec}{1 & 1 & 1 \\ 2 & 1 & 2 \\ 1 & 3 & 2}\end{myvec}
	\end{align}
	can be decomposed into the product of a lower triangular matrix $\vec{L}$ and an upper triangular matrix $\vec{U}$ as $\vec{A} = \vec{L}\vec{U}$. Let $\vec{x}, \vec{z} \in \mathbb{R}^3$ and $\vec{b} = (1 \;\; 1 \;\; 1)^\top$.
	\begin{enumerate}[label=\alph*)]
		\item The solution $\vec{z} = (z_1 \;\; z_2 \;\; z_3)^\top$ of the system $\vec{L}\vec{z} = \vec{b}$ is
		\begin{enumerate}[label=\roman*)]
			\item $(-1 \;\; -1 \;\; -2)^\top$
			\item $(1 \;\; -1 \;\; 2)^\top$
			\item $(1 \;\; -1 \;\; -2)^\top$
			\item $(-1 \;\; 1 \;\; 2)^\top$
		\end{enumerate}
		\item The solution $\vec{x} = (x_1 \;\; x_2 \;\; x_3)^\top$ of the system $\vec{U}\vec{x} = \vec{z}$ is
		\begin{enumerate}[label=\roman*)]
			\item $(2 \;\; 1 \;\; -2)^\top$
			\item $(2 \;\; 1 \;\; 2)^\top$
			\item $(-2 \;\; -1 \;\; -2)^\top$
			\item $(-2 \;\; 1 \;\; -2)^\top$
		\end{enumerate}
	\end{enumerate}
	
	\solution
	
	\textbf{LU Decomposition by Row Transformations}\\
	
	We will perform row operations on $\vec{A}$ to obtain $\vec{U}$. The multipliers used will form the entries of $\vec{L}$.
	\begin{align}
		\vec{A} = \begin{myvec}{1 & 1 & 1 \\ 2 & 1 & 2 \\ 1 & 3 & 2}\end{myvec}
	\end{align}
	Operation $R_2 \to R_2 - (\textcolor{red}{2})R_1$. The multiplier is $l_{21} = 2$.
	\begin{align}
		\begin{myvec}{1 & 1 & 1 \\ 0 & -1 & 0 \\ 1 & 3 & 2}\end{myvec}
	\end{align}
	Operation $R_3 \to R_3 - (\textcolor{red}{1})R_1$. The multiplier is $l_{31} = 1$.
	\begin{align}
		\begin{myvec}{1 & 1 & 1 \\ 0 & -1 & 0 \\ 0 & 2 & 1}\end{myvec}
	\end{align}
	Operation $R_3 \to R_3 - (\textcolor{red}{-2})R_2$. The multiplier is $l_{32} = -2$.
	\begin{align}
		\vec{U} = \begin{myvec}{1 & 1 & 1 \\ 0 & -1 & 0 \\ 0 & 0 & 1}\end{myvec}
	\end{align}
	The lower triangular matrix $\vec{L}$ is constructed from the multipliers.
	\begin{align}
		\vec{L} = \begin{myvec}{1 & 0 & 0 \\ l_{21} & 1 & 0 \\ l_{31} & l_{32} & 1}\end{myvec} = \begin{myvec}{1 & 0 & 0 \\ 2 & 1 & 0 \\ 1 & -2 & 1}\end{myvec}
	\end{align}
	
	\textbf{Part a): Solving $\vec{L}\vec{z} = \vec{b}$}\\
	
	We solve the system using Gaussian elimination on the augmented matrix $[\vec{L}|\vec{b}]$.
	\begin{align}
		\myaugvec{3}{1 & 0 & 0 & 1 \\ 2 & 1 & 0 & 1 \\ 1 & -2 & 1 & 1} 
		&\xleftrightarrow[R_3 \to R_3 - R_1]{R_2 \to R_2 - 2R_1} \myaugvec{3}{1 & 0 & 0 & 1 \\ 0 & 1 & 0 & -1 \\ 0 & -2 & 1 & 0} \\
		&\xleftrightarrow{R_3 \to R_3 + 2R_2} \myaugvec{3}{1 & 0 & 0 & 1 \\ 0 & 1 & 0 & -1 \\ 0 & 0 & 1 & -2}
	\end{align}
	From the reduced row echelon form, we get the solution $\vec{z} = (1 \;\; -1 \;\; -2)^\top$, which corresponds to option (iii).
	
	\textbf{Part b): Solving $\vec{U}\vec{x} = \vec{z}$}\\
	
	We solve the system using the augmented matrix $[\vec{U}|\vec{z}]$.
	\begin{align}
		\myaugvec{3}{1 & 1 & 1 & 1 \\ 0 & -1 & 0 & -1 \\ 0 & 0 & 1 & -2}
		&\xleftrightarrow{R_2 \to -R_2} \myaugvec{3}{1 & 1 & 1 & 1 \\ 0 & 1 & 0 & 1 \\ 0 & 0 & 1 & -2} \\
		&\xleftrightarrow[R_1 \to R_1 - R_2]{R_1 \to R_1 - R_3} \myaugvec{3}{1 & 0 & 0 & 2 \\ 0 & 1 & 0 & 1 \\ 0 & 0 & 1 & -2}
	\end{align}
	From the reduced row echelon form, we get the solution $\vec{x} = (2 \;\; 1 \;\; -2)^\top$, which corresponds to option (i).
	
	\begin{figure}[h!]
	\centering
	\includegraphics[width=0.8\columnwidth]{figs/Figure_1.png} 
	\caption*{Fig }
	\label{Fig1}
	\end{figure}
\end{document}